\begin{itemize}
	\item \textcolor{red}{The life-expectancy cost of kidney donation depends primarily on which published donor mortality hazard ratio (HR) is assumed: under two months at HR~$=1.0$ (no excess mortality) to over four years at the upper confidence bound of the longest-follow-up cohort. Under a time-varying base case that reconciles the existing cohort data (excess mortality emerging only after roughly a decade), the estimate is 2.4~years.}
	\item \textcolor{red}{This base-case cost is largest for the youngest donors. Black donors face a larger cost than White donors at young ages, but the difference disappears by age~55; men and women differ little at age~40, with the gap widening to 7--9\% at ages~25 and~55.}
	\item \textcolor{red}{Resolving the time-course of post-donation mortality risk is the dominant research priority: whether and when excess mortality emerges determines the dominant share of the life-expectancy cost.}
	\item \textcolor{red}{Priority waitlist access benefits individual donors by approximately 6~days at the population level, and designated non-donor family members (via a National Kidney Registry voucher) by under two days on average; either benefit rises to roughly 164--1{,}488~days if kidney failure actually occurs, depending on age at onset.}
\end{itemize}
